\documentclass[11pt]{article}

\usepackage[a4paper,margin=1in]{geometry}
\usepackage{amsmath,amssymb,amsthm,mathtools}
\usepackage{array}
\usepackage{booktabs}
\usepackage{enumitem}
\usepackage{tikz}
\usepackage{physics}
\usepackage{needspace}
\usepackage{xcolor}
\definecolor{courseblue}{RGB}{0,90,150}
\usetikzlibrary{arrows.meta,positioning}

\setlength{\parindent}{0pt}
\setlength{\parskip}{0.75em}

\newcounter{problem}
\newenvironment{problem}{%
    \refstepcounter{problem}%
    \Needspace{7\baselineskip}%
    \par\vspace{0.7em}%
    \noindent\textcolor{courseblue}{\rule{\textwidth}{0.4pt}}%
    \par\vspace{0.5em}%
    \begin{list}{}{%
        \setlength{\leftmargin}{2.3em}%
        \setlength{\labelwidth}{1.8em}%
        \setlength{\labelsep}{0.5em}%
        \setlength{\topsep}{0pt}%
        \setlength{\partopsep}{0pt}%
    }%
    \item[{[\theproblem]}]
}{%
    \end{list}
    \par\vspace{0.5em}
}

\newenvironment{parts}{%
    \begin{enumerate}[label={[\alph*]},leftmargin=2.2em,itemsep=0.65em,topsep=0.65em]
}{%
    \end{enumerate}
}

% Type each solution between \begin{answer} and \end{answer}.
\newenvironment{answer}{%
    \Needspace{4\baselineskip}%
    \par\vspace{0.35em}
    \noindent\textbf{Answer:}\par
}{%
    \par\vspace{0.25em}
}



\title{CO250 Assignment 4}
\author{}
\date{October 5--19}

\begin{document}
\maketitle

\begin{problem}
    The following LP is in SEF:
    \begin{equation}
    \begin{aligned}
        \max\quad & (1,-2,0,1,3)x \\
        \text{subject to}\quad &
        \begin{pmatrix}1&-1&2&-1&0\\2&0&1&-1&1\end{pmatrix}x
        =\begin{pmatrix}1\\-1\end{pmatrix} \\
        & x\geq\mathbf{0}
    \end{aligned}
    \end{equation}
    Find an equivalent LP in canonical form for:
    \begin{parts}
        \item the basis $\{1,4\}$,
        \begin{answer}
            \begin{align*}
                \max\quad & (1,-2,0,1,3)x \\
                \text{subject to}\quad &
                \begin{pmatrix}1&1&-1&0&1\\0&2&-3&1&1\end{pmatrix}x
                =\begin{pmatrix}-2\\-3\end{pmatrix} \\
                & x\geq\mathbf{0}
            \end{align*}
        \end{answer}
        \item the basis $\{3,5\}$.
        \begin{answer}
            \begin{align*}
                \max\quad & (1,-2,0,1,3)x \\
                \text{subject to}\quad &
                \begin{pmatrix}1/2&-1/2&1&-1/2&0\\3/2&1/2&0&-1/2&1\end{pmatrix}x
                =\begin{pmatrix}1/2\\-3/2\end{pmatrix} \\
                & x\geq\mathbf{0}
            \end{align*}
        \end{answer}
    \end{parts}
    In each case, state whether the basis (i.e.\ the corresponding basic feasible solution) is feasible.
\end{problem}

\begin{problem}
    Consider the following parameterized linear program in standard equality form, where $\alpha\in\mathbb{R}$:
    \begin{equation*}
    \begin{aligned}
        \max\quad & x_4 \\
        \text{s.t.}\quad & -x_1+x_3=1, \\
        & x_2+\alpha x_3+x_4=0, \\
        & x_1,x_2,x_3,x_4\geq 0.
    \end{aligned}
    \tag{$P_\alpha$}
    \end{equation*}
    \begin{parts}
        \item Find a certificate of infeasibility of $(P_\alpha)$ for $\alpha=1$.
        \begin{answer}

        \end{answer}
        \Needspace{7\baselineskip}\item Find a certificate of unboundedness of $(P_\alpha)$ for $\alpha=-1$.
        \begin{answer}

        \end{answer}
        \item Determine all values of $\alpha$ for which $(P_\alpha)$ is infeasible. Provide a certificate of infeasibility for each such $\alpha$.
        \begin{answer}

        \end{answer}
        \item Determine all values of $\alpha$ for which $(P_\alpha)$ is unbounded. Provide a certificate of unboundedness for each such $\alpha$.
        \begin{answer}

        \end{answer}
        \item For any remaining feasible values of $\alpha$, find an optimal solution $\bar{x}$ and a certificate for the optimality of $\bar{x}$.
        \begin{answer}

        \end{answer}
    \end{parts}
\end{problem}

\begin{problem}
    In this exercise, you will work through a guided application of the simplex method to an LP.
    You do not need to have learned all the details of the algorithm in order to complete the exercise.

    Consider the following LP
    \[
    \begin{aligned}
        \max\quad & (1\quad 1\quad 0\quad 0)x \\
        \text{s.t.}\quad &
        \begin{pmatrix}2&-1&1&0\\1&-2&0&1\end{pmatrix}x
        =\begin{pmatrix}1\\1\end{pmatrix} \\
        & x\geq\mathbf{0}.
    \end{aligned}
    \]
    \begin{parts}
        \item Let $\bar{x}$ be the basic feasible solution with basis $\{3,4\}$. What is $\bar{x}$?
        \begin{answer}

        \end{answer}
        \item Explain why $\bar{x}$ is not an optimal solution.
        \begin{answer}

        \end{answer}
        \Needspace{14\baselineskip}\item Let $x^{(\alpha)}$ be the solution to
        \[
            \begin{pmatrix}2&-1&1&0\\1&-2&0&1\end{pmatrix}x
            =\begin{pmatrix}1\\1\end{pmatrix}
        \]
        such that $x_1^{(\alpha)}=\alpha$ and $x_2^{(\alpha)}=0$.
        Determine the set
        \[
            \mathcal{A}=\{\alpha:x^{(\alpha)}\text{ is a feasible solution}\}.
        \]
        \begin{answer}

        \end{answer}
        \Needspace{9\baselineskip}\item Let $\alpha^*$ be the maximum value of $\alpha$ in $\mathcal{A}$ (from part (c) you should see that the maximum exists).
        Show that $x^{(\alpha^*)}$ is a basic feasible solution whose objective value is greater than that for $\bar{x}$.
        What is the basis for $x^{(\alpha^*)}$?
        \begin{answer}

        \end{answer}
        \Needspace{15\baselineskip}\item Redo part (c) but swapping the role of $x_1$ and $x_2$. Let $x^{(\alpha)}$ be the solution to
        \[
            \begin{pmatrix}2&-1&1&0\\1&-2&0&1\end{pmatrix}x
            =\begin{pmatrix}1\\1\end{pmatrix}
        \]
        such that $x_1^{(\alpha)}=0$ and $x_2^{(\alpha)}=\alpha$.
        Determine the set
        \[
            \mathcal{B}=\{\alpha:x^{(\alpha)}\text{ is a feasible solution}\}.
        \]
        What does $\mathcal{B}$ imply about the three possible outcomes of this LP?
        \begin{answer}

        \end{answer}
    \end{parts}
\end{problem}

\begin{problem}
    Consider the following linear program in standard equality form:
    \begin{equation*}
    \begin{aligned}
        \min\quad & x_4 \\
        \text{s.t.}\quad & x_1+x_2+x_3=3, \\
        & x_1+2x_2+3x_3+x_4=6, \\
        & x_1,x_2,x_3,x_4\geq 0.
    \end{aligned}
    \tag{P}
    \end{equation*}
    In matrix form, the constraints are $Ax=b$, $x\geq 0$, where
    \[
        A=\begin{pmatrix}1&1&1&0\\1&2&3&1\end{pmatrix},
        \qquad b=\begin{pmatrix}3\\6\end{pmatrix}.
    \]
    Let $\bar{x}$ be a feasible solution given by $(1\ 1\ 1\ 0)^T$.
    Let $J=\{1,2,3\}$ be the set of indices corresponding to the positive coordinates of $\bar{x}$.
    \begin{parts}
        \Needspace{10\baselineskip}\item Find a nonzero vector $d\in\mathbb{R}^4$ such that
        \[
            Ad=0
        \]
        and $d$ is supported on $J$. That is, $d_j=0$ for every $j\notin J$.
        \begin{answer}

        \end{answer}
        \Needspace{11\baselineskip}\item For $\alpha\in\mathbb{R}$, define
        \[
            x^{(\alpha)}=\bar{x}+\alpha d.
        \]
        Show that $Ax^{(\alpha)}=b$ for every $\alpha\in\mathbb{R}$.
        Determine $\mathcal{A}$, the set of all values of $\alpha$ for which $x^{(\alpha)}$ is feasible.
        \begin{answer}

        \end{answer}
        \Needspace{10\baselineskip}\item Choose $\alpha^*$ to be either $\max\{\alpha:\alpha\in\mathcal{A}\}$ or $\min\{\alpha:\alpha\in\mathcal{A}\}$, whichever exists.
        Compute the resulting feasible solution $x^*=x^{(\alpha^*)}$ and compare the number of its positive coordinates with that of $\bar{x}$.
        \begin{answer}

        \end{answer}
        \Needspace{15\baselineskip}\item Prove that every feasible LP in SEF has a basic feasible solution by generalising the above arguments as follows.
        Consider LP: $\max\{c^Tx:Ax=b,x\geq\mathbf{0}\}$.
        Let $\bar{x}$ be a feasible solution with minimum number of positive entries.
        Let $J$ be the set of indices corresponding to the positive coordinates of $\bar{x}$.
        Following the guidance below to show that $\{A_j:j\in J\}$ is linearly independent.
        \begin{enumerate}[label={(\roman*)},leftmargin=2.5em,itemsep=0.65em,topsep=0.65em]
            \item Suppose on the contrary that $\{A_j:j\in J\}$ is not linearly independent.
            Show that there exists a nonzero $d$ supported on $J$ such that $Ad=0$.
            \begin{answer}

            \end{answer}
            \item Explain why without loss of generality we may assume that $d_j<0$ for some $j\in J$.
            \begin{answer}

            \end{answer}
            \Needspace{10\baselineskip}\item For $\alpha\in\mathbb{R}$, define
            \[
                x^{(\alpha)}=\bar{x}+\alpha d.
            \]
            Show that $Ax^{(\alpha)}=b$ for every $\alpha\in\mathbb{R}$.
            \begin{answer}

            \end{answer}
            \item Let $\mathcal{A}$ be the set of all values of $\alpha$ for which $x^{(\alpha)}$ is feasible.
            Show that $\mathcal{A}$ has a maximum.
            \begin{answer}

            \end{answer}
            \Needspace{9\baselineskip}\item Let $\alpha^*$ be the maximum of $\mathcal{A}$ and let $x^*=x^{(\alpha^*)}$.
            Show that the number of positive coordinates in $x^*$ is smaller than $|J|$.
            \begin{answer}

            \end{answer}
            \item Conclude that $\{A_j:j\in J\}$ is linear independent.
            Explain why $\bar{x}$ is a basic feasible solution.
            \begin{answer}

            \end{answer}
        \end{enumerate}
    \end{parts}
\end{problem}

\end{document}
